\begin{frame}
	\frametitle{¿Por qué este curso?}

	Muchos fenómenos en la naturaleza se describen mediante ecuaciones
	diferenciales parciales (EDPs).

	\begin{columns}
		\begin{column}{.5\paperwidth}
			\begin{block}{Fenómenos físicos}
				\begin{description}
					\item[Propagación de ondas]

					      tsunamis, terremotos, sonido.

					\item[Difusión térmica]

					      distribución de calor en materiales.

					\item[Dinámica de fluidos]

					      océanos, atmósfera, aire en tuberías.

					\item[Reacciones químicas]

					      combustión, ecosistemas.
				\end{description}
			\end{block}

			La mayoría de EDPs no tienen soluciones analíticas.

			\begin{center}\large
				¡Necesitamos métodos numéricos para resolverlas!
			\end{center}
		\end{column}
		\begin{column}{.4\paperwidth}
			\begin{figure}[ht!]
				\centering
				\includegraphics[width=0.4\paperwidth]{geoclaw}
				\caption{Simulación numérica de un tsunami (Fuente: \url{https://dlgeorge.github.io/project/geoclaw-project})}
				\label{fig:tsunami}
			\end{figure}
		\end{column}
	\end{columns}
\end{frame}

\begin{frame}
	\frametitle{Objetivos del curso}

	\begin{columns}
		\begin{column}{.45\paperwidth}
			Al final de este curso, ustedes serán capaces de

			\begin{enumerate}
				\item \textbf{Comprender} la teoría detrás de los métodos numéricos.
				      \begin{itemize}
					      \item Consistencia, convergencia, estabilidad.
					      \item Análisis de errores.
				      \end{itemize}

				\item \textbf{Implementar} esquemas numéricos en Python.
				      \begin{itemize}
					      \item Métodos de diferencias finitas (FDM).
					      \item Métodos de elementos finitos (FEM).
					      \item Métodos de volúmenes finitos (FVM).
				      \end{itemize}

				\item \textbf{Resolver} problemas prácticos.
				      \begin{itemize}
					      \item Ecuaciones parabólicas (calor).
					      \item Ecuaciones elípticas (Poisson).
					      \item Ecuaciones hiperbólicas (ondas, leyes de conservación).
				      \end{itemize}
			\end{enumerate}
		\end{column}
		\begin{column}{.35\paperwidth}
			\begin{figure}[ht!]
				\centering
				\includegraphics[width=0.4\paperwidth]{githubclassroom}
				\caption{Simulación numérica de la ecuación de calor (Fuente: Wikipedia)}
				\label{fig:calor}
			\end{figure}
		\end{column}
	\end{columns}
\end{frame}

\begin{frame}
	\frametitle{Dinámica del curso}

	\begin{block}{6 horas por semana durante 16 semanas}
		\begin{table}[ht!]
			\centering
			\begin{tabular}{|l|l|}
				\hline
				\textbf{Teoría}      & 2h Lunes + 2h Miércoles      \\[0.1cm]
				\hline
				\textbf{Laboratorio} & 2h Martes (Python + Jupyter) \\[0.1cm]
				\hline
			\end{tabular}
		\end{table}
	\end{block}

	\begin{block}{Evaluación continua}
		\begin{itemize}
			\item 5 prácticas calificadas (semanas 3, 5, 10, 12, 15).
			\item 1 examen parcial (semana 8).
			\item 1 proyecto con 2 presentaciones (semanas 7 y 15).
			\item Trabajo de laboratorio regular.
		\end{itemize}
		\small Énfasis en argumentación clara y rigor matemático.
	\end{block}
\end{frame}

\begin{frame}
	\frametitle{GitHub Classroom (\url{https://classroom.github.com})}

	\begin{columns}
		\begin{column}{.45\paperwidth}
			\begin{block}{Plataforma de trabajo colaborativo}
				\begin{itemize}
					\item \textbf{Repositorio público} para cada estudiante que permite mantener un historial completo de cambios y versiones.
					\item \textbf{Control de versiones} con Git.
					\item \textbf{Automatización} de retroalimentación.
					\item Acceso a \textbf{recursos compartidos} de la clase.
				\end{itemize}
			\end{block}

			\begin{block}{Ventajas}
				\begin{itemize}
					\item Aprenderán herramientas profesionales de desarrollo.
					\item Tendrán registro completo de su trabajo.
					\item Facilita la evaluación y retroalimentación.
					\item Prepara para colaboración en proyectos reales.
				\end{itemize}
			\end{block}
		\end{column}
		\begin{column}{.35\paperwidth}
			\begin{figure}[ht!]
				\centering
				\includegraphics[width=0.4\paperwidth]{githubdesktop}
				\caption{Interfaz de GitHub Classroom (Fuente: GitHub)}
				\label{fig:github}
			\end{figure}
		\end{column}
	\end{columns}
\end{frame}

\begin{frame}
	\frametitle{Contenido: Ruta de aprendizaje}

	\begin{block}{Bloque 1: Fundamentos (Semanas 1-3)}
		\begin{itemize}
			\item Introducción al método de las diferencias finitas.
			\item Consistencia, convergencia, estabilidad.
			\item Método de Galerkin.
		\end{itemize}
	\end{block}

	\begin{block}{Bloque 2: Ecuaciones hiperbólicas (Semanas 4-7)}
		\begin{itemize}
			\item Leyes de conservación y ondas de choque.
			\item Esquemas clásicos (Upwind, Lax-Friedrichs, Lax-Wendroff).
			\item Aplicación: ecuaciones de agua poco profundas.
		\end{itemize}
	\end{block}

	\begin{block}{Bloque 3: Métodos sofisticados (Semanas 9-15)}
		\begin{itemize}
			\item Métodos de volúmenes finitos (FVM).
			\item Métodos WENO/ENO.
			\item Esquemas que preservan la variación total (TVD).
			\item Aplicaciones a problemas reales.
		\end{itemize}
	\end{block}
\end{frame}

\begin{frame}
	\frametitle{Recursos del curso}

	\begin{block}{Sitio web (\url{https://cm032.github.io})}
		Encontrarás materiales y referencias como
		\begin{itemize}
			\item Diapositivas de clase.
			\item Notas de laboratorio.
			\item Ejercicios resueltos.
			\item Lecturas complementarias.
			\item Cronograma detallado del curso (fechas de entregas, temas, etc.):  \url{https://cm032.github.io/assets/misc/cronogram.pdf}
		\end{itemize}
	\end{block}

	\begin{block}{Libros de referencia principales}
		\begin{itemize}
			\item LeVeque (2007) — Finite Difference Methods for ODEs and PDEs
			\item Thomas (1995) — Numerical Partial Differential Equations
			\item Vázquez-Cendón — Finite Volume Methods for Hyperbolic Problems
		\end{itemize}
	\end{block}

	\begin{block}{Herramientas}
		\begin{itemize}
			\item Python 3, Jupyter Notebooks
			\item Git y GitHub Classroom
		\end{itemize}
	\end{block}
\end{frame}

\begin{frame}
	\frametitle{¿Estamos listos?}

	\centering
	\vspace{1cm}

	{\Large \textbf{Vamos a aprender juntos cómo}}

	{\Large \textbf{capturar el comportamiento de la naturaleza}}

	{\Large \textbf{mediante ecuaciones numéricas.}}

	\vspace{1cm}

	{\large Preguntas y dudas: No duden en preguntar en cualquier momento}

	\vspace{1cm}

	¡Bienvenidos al CM032!
\end{frame}
